\RequirePackage{expl3}
\ExplSyntaxOn
\cs_generate_variant:Nn \tl_if_eq:nnT { VnT }
\cs_generate_variant:Nn \tl_if_eq:nnF { VnF }
\cs_generate_variant:Nn \tl_if_eq:nnTF { VnTF }
\ExplSyntaxOff

\documentclass[12pt,twoside=false,a4paper,parskip]{scrbook}
% use the following line instead of the class above if you want double-sided printing
% don't remove openright as it looks better in case of double-sided print
% \documentclass[12pt,twoside,a4paper,parskip,openright]{scrbook} 
\usepackage[utf8]{inputenc}
\usepackage[T1]{fontenc}
\usepackage{lmodern}
\usepackage{microtype}
\usepackage{csquotes}
\usepackage[ngerman]{babel}
\usepackage{subcaption} % modern alternative to subfigure
\usepackage[pdftex]{graphicx}
\usepackage[hidelinks]{hyperref}
\usepackage{xcolor}
\usepackage{amssymb}
\usepackage{textcomp}
\usepackage{nicefrac}
\usepackage{pdfpages}
\usepackage{float}
\usepackage{pdflscape}
\usepackage[verbose]{placeins}
\usepackage[markcase=ignoreuppercase,headsepline,plainfootsepline]{scrlayer-scrpage}
\usepackage{listings}
\usepackage{caption}
\usepackage{longtable}
\usepackage{setspace}
\usepackage{booktabs}
\usepackage[style=ieee,backend=biber,sorting=none]{biblatex}
\bibliography{literatur}
\usepackage[htt]{hyphenat} % texttt do line breaks
\usepackage{tabularx}
\usepackage{multirow}
\usepackage{enumitem}


%%%%%%%%%%%%%%%%%%%
%% definitions
%%%%%%%%%%%%%%%%%%%
\def\BaAuthor{Martin Rapps}
\def\BaAuthorStudyProgram{Geovisualisierung} 
\def\BaType{Projektarbeit} 
\def\BaTitle{KI-gestützte 3D-Rekonstruktion linearer Infrastruktur: Evaluierung einer Docker-basierten Scan-to-BIM Pipeline mittels SAM 3 und Gaussian Splatting}
\def\BaSupervisorOne{Prof. Dr. Müller}
\def\BaDeadline{[Abgabedatum]}

\def\iswithfullname{1} % uncomment this to create non-anonymous version
\ifdefined\iswithfullname
  \def\ShowBaAuthor{\BaAuthor}
\else
  \def\ShowBaAuthor{N.~N.}
\fi

\hypersetup{
pdfauthor={\ShowBaAuthor},
pdftitle={\BaTitle},
pdfsubject={Projektarbeit},
pdfkeywords={3DGS, Scan-to-BIM, SAM 3, SuGaR, Docker, Georeferenzierung}
}

%%%%%%%%%%%%%%%%%%%
%% configs to include
%%%%%%%%%%%%%%%%%%%
\colorlet{punct}{red!60!black}
\definecolor{background}{HTML}{EEEEEE}
\definecolor{delim}{RGB}{20,105,176}
\colorlet{numb}{magenta!60!black}

\definecolor{gray}{rgb}{0.4,0.4,0.4}
\definecolor{darkblue}{rgb}{0.0,0.0,0.6}
\definecolor{cyan}{rgb}{0.0,0.6,0.6}

\definecolor{pblue}{rgb}{0.13,0.13,1}
\definecolor{pgreen}{rgb}{0,0.5,0}
\definecolor{pred}{rgb}{0.9,0,0}
\definecolor{pgrey}{rgb}{0.46,0.45,0.48}

\lstset{
  basicstyle=\ttfamily,
  columns=fullflexible,
  showstringspaces=false,
  commentstyle=\color{gray}\upshape,
  linewidth=\textwidth
}

\lstdefinelanguage{json}{
    basicstyle=\normalfont\ttfamily,
    numbers=left,
    numberstyle=\scriptsize,
    stepnumber=1,
    numbersep=8pt,
    showstringspaces=false,
    breaklines=true,
    backgroundcolor=\color{background},
    literate=
     *{0}{{{\color{numb}0}}}{1}
      {1}{{{\color{numb}1}}}{1}
      {2}{{{\color{numb}2}}}{1}
      {3}{{{\color{numb}3}}}{1}
      {4}{{{\color{numb}4}}}{1}
      {5}{{{\color{numb}5}}}{1}
      {6}{{{\color{numb}6}}}{1}
      {7}{{{\color{numb}7}}}{1}
      {8}{{{\color{numb}8}}}{1}
      {9}{{{\color{numb}9}}}{1}
      {:}{{{\color{punct}{:}}}}{1}
      {,}{{{\color{punct}{,}}}}{1}
      {\{}{{{\color{delim}{\{}}}}{1}
      {\}}{{{\color{delim}{\}}}}}{1}
      {[}{{{\color{delim}{[}}}}{1}
      {]}{{{\color{delim}{]}}}}{1},
}

\lstset{language=xml,
  morestring=[b]",
  morestring=[s]{>}{<},
  morecomment=[s]{<?}{?>},
  stringstyle=\color{black},
  numbers=left,
  numberstyle=\scriptsize,
  stepnumber=1,
  numbersep=8pt,
  identifierstyle=\color{darkblue},
  keywordstyle=\color{cyan},
  backgroundcolor=\color{background},
  morekeywords={xmlns,version,type}% list your attributes here
}

\newcommand*{\forcetwosidetitle}[1][1]{%
 \begingroup
   \cleardoubleoddpage
   \KOMAoptions{titlepage=true}% useful e.g. for scrartcl
   \csname @twosidetrue\endcsname
   \maketitle[{#1}]
 \endgroup
}


\begin{document}


%%%%%%%%%%%%%%%%%%%
%% Titelseite
%%%%%%%%%%%%%%%%%%%


\frontmatter
\titlehead{
  {Technische Hochschule Würzburg-Schweinfurt\\
    Fakultät Kunststofftechnik und Vermessung}}
\subject{\BaType}
\title{\BaTitle\\[15mm]}
\subtitle{\normalsize{vorgelegt an der Technischen Hochschule W\"{u}rzburg-Schweinfurt in der Fakult\"{a}t Kunststofftechnik und Vermessung zum Abschluss des Moduls Wissenschaftliches Arbeiten im Studiengang \BaAuthorStudyProgram}}
\author{\ShowBaAuthor}
\date{\normalsize{Eingereicht am: \BaDeadline}}
\publishers{
  \normalsize{Erstprüfer: \BaSupervisorOne}\\
}

\forcetwosidetitle


%%%%%%%%%%%%%%%%%%%
%% abstract
%%%%%%%%%%%%%%%%%%%

\addchap{Kurzfassung / Abstract}

[Hier kommt eine kurze Zusammenfassung (ca. halbe bis ganze Seite) der Motivation, Methodik und wichtigsten Ergebnisse deiner Projektarbeit rein.]


%%%%%%%%%%%%%%%%%%%
%% Inhaltsverzeichnis
%%%%%%%%%%%%%%%%%%%
\tableofcontents


%%%%%%%%%%%%%%%%%%%
%% Main part of the thesis
%%%%%%%%%%%%%%%%%%%
\mainmatter


\chapter{Einleitung und Problemstellung} \label{sec:einleitung}

\section{Motivation und Praxisbezug}
% Die Relevanz linearer Infrastrukturen (TenneT Anwendungsfall) und die Notwendigkeit von Scan-to-BIM.

\section{Forschungsfragen der Projektarbeit}
% Fokus der PA: 3DGS vs. klassische Photogrammetrie bezüglich detailarmer und feiner Strukturen.

\section{Struktur der Arbeit}
% Aufbau dieser Ausarbeitung.


\chapter{Grundlagen und Stand der Technik} \label{sec:grundlagen}

\section{Structure-from-Motion (SfM) vs. Multi-View Stereo (MVS)}
% Grundlagen der photogrammetrischen Rekonstruktion mittels COLMAP.

\section{3D Gaussian Splatting (3DGS)}
% Mathematische Grundlagen des Splattings nach Kerbl et al. (2023).

\section{Kombination mit 2D-Segmentierungsmodellen (SAM 3)}
% Ansätze zur Erhöhung der Geometrieauflösung durch vorgelagerte Maskierung.


\chapter{Konzeption der Docker-Infrastruktur} \label{sec:konzeption}

\section{Microservice-Architektur der Pipeline}
% Vorstellung der 5 Container (SAM 3, COLMAP, STS, SuGaR, Post-Processing).

\section{Vermeidung von Versionierungskonflikten}
% CUDA und PyTorch Versionen in den einzelnen Containern (Bind Mounts, WSL2).

\section{Automatisierte Orchestrierung und Breakpoint-Design}
% One-Click-Orchestration-Skript und der manuelle Schritt zur Passpunkt-Identifikation.


\chapter{Die Georeferenzierungs-Pipeline} \label{sec:georef}

\section{Das Georeferenzierungs-Paradoxon}
% Warum die Ausblendung der Umgebung durch STS die GCP-Zuordnung nach dem Training unmöglich macht.

\section{Pre-Training Georeferenzierung}
% Relative Koordinaten vs. Vertex-Jittering bei großen UTM-Koordinaten.

\section{Transformations-Matrix-Berechnung}
% Helmert-7-Parameter-Transformation via CloudCompare (Point-Picking).


\chapter{Umsetzung und Praktische Erprobung} \label{sec:umsetzung}

\section{Testdatenaufbau und Bilderfassung}
% Aufnahmekonfiguration der Aluabsperrung als PoC-Objekt.

\section{Segment-then-Splat Training}
% Maskierung mit SAM 3 und Training des objektspezifischen Loss.

\section{SuGaR Surface-Extraction}
% Geometrische Regularisierung und Erzeugung des Kabel-Sub-Meshes.

\section{Post-Processing in Container E}
% DGtal-Centerline und Transformation ins globale System.


\chapter{Ergebnisse, Evaluation und Diskussion} \label{sec:ergebnisse}

\section{Photometrischer Vergleich (PSNR, SSIM, LPIPS)}
% Optische Qualität der Novel View Synthesis.

\section{Geometrischer Vergleich: MVS vs. 3DGS}
% Rekonstruktionstiefe an feinen Strukturen der Aluabsperrung.

\section{Benchmark der Rechenzeiten}
% Vergleich der dichten MVS-Rekonstruktion vs. 3DGS-Trainingszeit.

\section{Diskussion der Limitierungen}
% VRAM-Bedarf, Skriptstabilität und Docker-I/O unter WSL2.


\chapter{Zusammenfassung und Ausblick auf die Bachelorarbeit} \label{sec:fazit}

\section{Zusammenfassendes Fazit}
% Beantwortung der Forschungsfragen der PA.

\section{Ausblick und Erweiterungen für die Bachelorarbeit}
% Geplante Feldarbeit, reale GNSS-Kabelmessung und RMSE/Hausdorff-Evaluation.


\backmatter
%%%%%%%%%%%%%%%%%%%
%% create figure list
%%%%%%%%%%%%%%%%%%%

\listoffigures
\addcontentsline{toc}{chapter}{Verzeichnisse}

%%%%%%%%%%%%%%%%%%%
%% create appendix
%%%%%%%%%%%%%%%%%%%


\appendix
\addchap{Anhänge}

\addsec{Anhang A: Installations- und Ausführungsskripte} \label{appendix:skripte}
% Master-Skript und Dockerfiles

\addsec{Anhang B: Deklaration über den Einsatz von Künstlicher Intelligenz} \label{appendix:ki_deklaration}
\begin{itemize}
  \item \textbf{Verwendetes KI-System:} [z.B. Gemini 3.5 Flash]
  \item \textbf{Verwendete Entwicklungsumgebung:} [z.B. Antigravity IDE]
  \item \textbf{Datum der Nutzung:} \today
  \item \textbf{Wesentliche Prompts (Auszug):}
        \begin{itemize}
          \item \enquote{Erstelle mir einen LaTeX-Entwurf basierend auf meinem Exposé Expose\_PA\_BA.}
        \end{itemize}
\end{itemize}

%%%%%%%%%%%%%%%%%%%
%% create tables list
%%%%%%%%%%%%%%%%%%%
\listoftables

%%%%%%%%%%%%%%%%%%%
%% create listings list
%%%%%%%%%%%%%%%%%%%
%\lstlistoflistings
%\addcontentsline{toc}{chapter}{Listings}

\cleardoublepage
\phantomsection
\addcontentsline{toc}{chapter}{Literatur}
\printbibliography

%%%%%%%%%%%%%%%%%%%
%% declaration on oath
%%%%%%%%%%%%%%%%%%%

\addchap{Eidesstattliche Erklärung}

Hiermit versichere ich, dass ich die vorliegende Arbeit eigenständig verfasst und keine anderen als die angegebenen Quellen und Hilfsmittel verwendet habe.

Alle übernommenen Inhalte sowie mit Unterstützung von KI generierten Inhalte wurden entsprechend den anerkannten wissenschaftlichen Grundsätzen oder entsprechend der Regelungen zur Kennzeichnung von KI-Inhalten kenntlich gemacht. Ausgenommen von der Kenntlichmachung sind orthografische oder grammatikalische Korrekturen, Übersetzungen sowie nicht-sinnverändernde Verbesserungen von Formulierungen. Ich bin mir bewusst, dass mit KI generierte Texte keine Garantie für die Qualität von Inhalten und Text bieten. Daher erkläre ich, dass ich KI-Werkzeuge lediglich als Hilfsmittel genutzt habe, die von KI generierten Inhalte kritisch überprüft habe und mein eigenständiger kognitiver sowie kreativer Einfluss in dieser Arbeit überwiegt. Ich versichere, dass ich die Inhalte meiner Arbeit vollständig verstanden habe und selbstständig vertreten kann.

\vspace{20pt}
\begin{flushright}
  $\overline{~~~~~~~~~~~~~~~~~\mbox{\ShowBaAuthor, am \today}~~~~~~~~~~~~~~~~~}$
\end{flushright}

\addchap{Zustimmung zur Plagiatsüberprüfung}

Hiermit willige ich ein, dass zum Zwecke der Überprüfung auf Plagiate meine vorgelegte Arbeit in digitaler Form an PlagScan (www.plagscan.com) übermittelt und diese vorübergehend (max. 5~Jahre) in der von PlagScan geführten Datenbank gespeichert wird sowie persönliche Daten, die Teil dieser Arbeit sind, dort hinterlegt werden.

\begin{small}
  Die Einwilligung ist freiwillig. Ohne diese Einwilligung kann unter Entfernung aller persönlichen Angaben und Wahrung der urheberrechtlichen Vorgaben die Plagiatsüberprüfung nicht verhindert werden. Die Einwilligung zur Speicherung und Verwendung der persönlichen Daten kann jederzeit durch Erklärung gegenüber der Fakultät widerrufen werden.
\end{small}

\vspace{20pt}
\begin{flushright}
  $\overline{~~~~~~~~~~~~~~~~~\mbox{\ShowBaAuthor, am \today}~~~~~~~~~~~~~~~~~}$
\end{flushright}

\end{document}
